\section*{Глава 1. Анализ предметной области}
\addcontentsline{toc}{section}{Глава 1. Анализ предметной области}
\stepcounter{section}

\subsection{Анализ особенностей работы с конфигурациями и диагностическими данными ПЛК LogicBox в системах АСУТП}

Автоматизированная система управления технологическим процессом (АСУТП)
представляет собой совокупность программных и технических средств,
обеспечивающих сбор данных, обработку сигналов, формирование управляющих
воздействий и контроль состояния технологического оборудования. Одним из
ключевых элементов такой системы является программируемый логический
контроллер (ПЛК), на котором исполняются алгоритмы управления, хранятся
параметры и настройки, а также формируются диагностические сведения о состоянии
устройства и подключенных подсистем. В связи с этим сопровождение ПЛК не может
рассматриваться как работа с отдельным файлом проекта: фактически речь идет о
работе с совокупностью конфигураций, журналов, резервных копий, служебных
комментариев и сведений об объекте эксплуатации.

Для контроллеров LogicBox практическая задача осложняется тем, что инженерное
взаимодействие с устройствами обычно происходит в локальной сети объекта по
Ethernet. Открытый репозиторий LogicBox содержит host-side software для
конфигурации и контроля ПЛК, разработанное под Debian GNU/Linux; в нем
присутствуют утилиты \texttt{lbdiscover}, \texttt{lbcmd}, \texttt{lbconf},
\texttt{lbgetlog}, \texttt{lbota}, \texttt{lbupload}, а для работы требуется
компьютер в той же Ethernet LAN, что и базовый модуль LogicBox~\cite{logicbox_repo}.
Следовательно, в рассматриваемой предметной области уже существует низкоуровневый
инструментальный слой, однако отсутствует единая веб-ориентированная среда,
которая позволяла бы хранить результаты работы с этими утилитами в общей базе,
разграничивать доступ пользователей и связывать технические данные с
эксплуатационным контекстом.

В инженерной практике предметная область включает следующие типы данных:
\begin{itemize}
    \item сведения об объектах автоматизации, площадках, шкафах и конкретных ПЛК;
    \item конфигурационные данные контроллеров, версии параметров и снимки состояний;
    \item диагностические записи: ошибки, предупреждения, сервисные события, журналы;
    \item файловые вложения: архивы, выгрузки, резервные копии, отчеты, пояснения;
    \item информация об источнике данных: кто, когда и на каком объекте получил эти сведения.
\end{itemize}

Особенностью работы в АСУТП является то, что инженер не всегда имеет доступ к
сети Интернет, но часто имеет доступ к локальной сети объекта, в которой
находятся ПЛК. Поэтому для данной темы целесообразно рассматривать веб-систему
не как средство прямого удаленного управления оборудованием, а как центральный
контур учета и анализа данных. Получение информации из локальной сети может
выполняться вспомогательным локальным клиентом, который обнаруживает
контроллеры, считывает сведения штатными утилитами и передает результаты на
сервер при наличии сетевого соединения. В рамках курсового проекта именно
веб-приложение является основным объектом разработки, а локальный клиент
рассматривается как интеграционный внешний компонент.

Такой подход позволяет сохранить акцент на дисциплине <<Разработка
веб-приложений>> и одновременно сделать проект практически полезным. Центральный
сервер обеспечивает хранение данных, поиск, фильтрацию, управление доступом и
аудит действий пользователей, а локальный клиент выполняет только сбор и
передачу данных. Исключение функций удаленного управления ПЛК из первой версии
системы дополнительно снижает риски, связанные с безопасностью промышленной
инфраструктуры.

На основании анализа предметной области можно выделить основные проблемы,
которые должна решать разрабатываемая система:
\begin{enumerate}
    \item отсутствие единого веб-реестра объектов эксплуатации, контроллеров и связанных конфигураций;
    \item сложность определения актуальной версии конфигурации и источника ее получения;
    \item разрыв между диагностическими журналами и эксплуатационным контекстом объекта;
    \item отсутствие удобного механизма накопления и последующего анализа данных, собранных инженером в локальной сети;
    \item необходимость строгого разграничения доступа, при котором пользователь видит только разрешенные ему объекты и ПЛК.
\end{enumerate}

Сформулированные особенности предметной области приведены к требованиям,
представленным в таблице~\ref{tab:domain_requirements}.

\begin{table}[ht!]
    \centering
    \small
    \caption{Ключевые требования предметной области}
    \label{tab:domain_requirements}
    \begin{tabular}{|p{0.22\textwidth}|p{0.30\textwidth}|p{0.35\textwidth}|}
        \hline
        \textbf{Область} & \textbf{Практическая потребность} & \textbf{Следствие для веб-приложения} \\ \hline
        Учет ПЛК и объектов & Видеть, какие контроллеры относятся к конкретной площадке, шкафу и АСУТП & Ведение структурированного реестра объектов, узлов и контроллеров \\ \hline
        Конфигурации & Хранить актуальные и архивные версии данных, понимать источник получения и дату фиксации & Карточки конфигураций, история версий, атрибуты автора и времени загрузки \\ \hline
        Диагностика & Быстро находить причины повторяющихся сбоев и анализировать журналы событий & Централизованное хранение диагностических записей, поиск, фильтрация и агрегирование \\ \hline
        Работа без Интернета & Сохранять результаты локальной инженерной работы и передавать их позже & Поддержка приема данных от внешнего локального клиента и учета сессий синхронизации \\ \hline
        Безопасность доступа & Ограничивать обзор и изменение данных по роли и области ответственности & Ролевая модель, привязка прав к объектам и аудит действий пользователей \\ \hline
    \end{tabular}
    \normalsize
\end{table}

Таким образом, предметная область требует создания не универсального файлового
архива, а специализированного веб-приложения, ориентированного на учет ПЛК,
связанных конфигураций, диагностических данных и истории взаимодействия с
локальными инженерными инструментами.

\subsection{Обзор существующих программных решений для учета конфигураций, журналов событий и сопровождения PLC-проектов}

Существующие решения целесообразно рассматривать по группам задач. Одни системы
ориентированы на управление инженерными проектами, другие --- на хранение версий,
третьи --- на сбор журналов, а четвертые представляют собой низкоуровневые
утилиты для взаимодействия с устройствами. Для темы данного курсового проекта
важно показать, что требуемое веб-приложение не дублирует полностью уже
существующие продукты, а объединяет их сильные стороны применительно к
контроллерам LogicBox и внутренним процессам сопровождения.

CODESYS Automation Server представляет собой зрелое решение для централизованной
работы с устройствами и приложениями в экосистеме CODESYS. Его сильной стороной
являются развитые механизмы учета устройств, удаленного обслуживания,
управления версиями и правами пользователей~\cite{codesys_as}. Однако данная
платформа ориентирована на собственную технологическую среду и не предназначена
для легковесной интеграции именно с LogicBox в формате внутреннего
специализированного веб-приложения.

Siemens TIA Portal Project Server решает задачу коллективной работы с
инженерными проектами и обеспечивает распределение ролей при разработке и
сопровождении проектов автоматизации~\cite{siemens_multiuser}. Для крупных
проектов это сильный подход, но он также привязан к экосистеме Siemens и в
первую очередь ориентирован на совместную инженерную разработку, а не на
формирование независимого корпоративного веб-реестра оборудования и диагностики.

GitLab полезен как платформа для хранения версий, управления задачами и
организации командной разработки~\cite{gitlab_platform,gitlab_issues}. Он
хорошо решает вопросы истории изменений, комментариев и контроля доступа к
репозиториям, но не содержит предметной модели ПЛК, не умеет из коробки
учитывать объекты эксплуатации и не решает задачу структурированного хранения
снимков диагностического состояния оборудования.

Grafana Loki предназначена для централизованного сбора и поиска логов,
позволяет агрегировать события и строить аналитические выборки~\cite{loki_overview}.
Это делает ее полезной с точки зрения анализа журналов, однако Loki не является
готовой системой сопровождения ПЛК, поскольку не связывает журналы с карточками
контроллеров, конфигурациями, объектами эксплуатации и ролевыми сценариями
работы инженеров.

Наконец, открытый репозиторий LogicBox представляет особый интерес для темы
проекта, поскольку уже содержит практические утилиты взаимодействия с ПЛК,
включая обнаружение устройств в сети и получение части конфигурационных и
диагностических данных~\cite{logicbox_repo}. Вместе с тем эти средства работают
на уровне командной строки и не предоставляют веб-интерфейс, общую базу данных,
управление ролями, механизм согласованной загрузки файлов и историю работы по
объектам. Следовательно, их целесообразно рассматривать не как конкурентное
решение, а как интеграционную основу для внешнего локального клиента, который
будет взаимодействовать с центральным веб-приложением.

Обобщенный сравнительный анализ приведен в таблице~\ref{tab:solutions_compare}.

\small
\begin{longtable}{|p{0.18\textwidth}|p{0.22\textwidth}|p{0.24\textwidth}|p{0.24\textwidth}|}
    \caption{Сравнительный анализ существующих решений} \label{tab:solutions_compare} \\
    \hline
    \textbf{Решение} & \textbf{Сильные стороны} & \textbf{Что закрывает хорошо} & \textbf{Ограничения для темы проекта} \\ \hline
    \endfirsthead

    \tabcont{\ref{tab:solutions_compare}}
    \textbf{Решение} & \textbf{Сильные стороны} & \textbf{Что закрывает хорошо} & \textbf{Ограничения для темы проекта} \\ \hline
    \endhead

    \hline
    \endfoot

    \hline
    \endlastfoot

    CODESYS Automation Server & Централизованная работа с устройствами, версиями и пользователями & Управление жизненным циклом устройств и проектов в экосистеме CODESYS & Привязка к собственной платформе, избыточность для узкоспециализированной внутренней системы \\ \hline
    Siemens TIA Portal Project Server & Развитая совместная работа с инженерными проектами и распределение ролей & Командная разработка и сопровождение проектов Siemens & Ориентация на TIA Portal, а не на независимое веб-приложение для LogicBox \\ \hline
    GitLab & Хранение версий, история изменений, задачи, self-hosted-развертывание & Версионирование файлов и организацию командной работы & Отсутствие предметной модели ПЛК, диагностических сущностей и интеграции с локальной сетью объекта \\ \hline
    Grafana Loki & Централизованный сбор логов, поиск и агрегирование событий & Анализ журналов и эксплуатационных сообщений & Не решает учет ПЛК, конфигураций, файловых вложений и ролей по объектам \\ \hline
    Утилиты LogicBox & Обнаружение устройств, чтение части конфигурационных и диагностических данных, работа в локальной сети объекта & Низкоуровневое взаимодействие с ПЛК LogicBox & Нет единого веб-интерфейса, общей БД, ролевой модели и истории сопровождения \\ \hline
\end{longtable}
\normalsize

Из проведенного анализа следует, что ни одно из рассмотренных средств не
закрывает в полном объеме задачу, поставленную в рамках курсового проекта.
Наиболее перспективной является модель, в которой центральное место занимает
веб-приложение, а существующие утилиты LogicBox используются косвенно --- через
внешний локальный клиент, отвечающий за сбор данных в сети объекта.

\subsection{Формулировка цели, задач и функциональных требований к разрабатываемому веб-приложению}

С учетом проведенного анализа цель курсового проекта может быть уточнена
следующим образом: разработать веб-приложение для централизованного учета
контроллеров LogicBox, хранения конфигурационных и диагностических данных,
поддержки ролевого доступа к объектам эксплуатации и приема результатов,
полученных локальным клиентом в сети объекта, без реализации функций удаленного
управления оборудованием.

Такое уточнение не противоречит уже утвержденной теме работы, акцент на веб-приложении полностью
сохраняется, а локальный клиент выступает только как вспомогательный источник
данных для серверной части системы.

Практическое назначение разрабатываемой системы заключается в следующем:
\begin{itemize}
    \item вести единый веб-реестр объектов эксплуатации, шкафов и контроллеров LogicBox;
    \item хранить конфигурации, снимки параметров, журналы и связанные вложения в структурированном виде;
    \item обеспечивать поиск, фильтрацию и агрегирование диагностических данных для анализа эксплуатации;
    \item принимать результаты локального сбора данных, выполненного инженером в сети объекта;
    \item ограничивать доступ к данным в соответствии с ролью пользователя и назначенной ему областью ответственности.
\end{itemize}

Для достижения поставленной цели необходимо решить следующие задачи:
\begin{enumerate}
    \item проанализировать особенности сопровождения конфигураций и диагностических данных ПЛК LogicBox в АСУТП;
    \item определить состав сущностей, подлежащих учету во внутренней предметной модели веб-приложения;
    \item разработать ролевую модель и правила разграничения доступа к объектам эксплуатации;
    \item спроектировать веб-интерфейсы для просмотра, создания, редактирования и удаления основных сущностей;
    \item реализовать загрузку, скачивание и хранение конфигурационных файлов, журналов и сопроводительных документов;
    \item реализовать серверный механизм приема данных от внешнего локального клиента;
    \item реализовать поиск, фильтрацию и агрегирование накопленных диагностических данных.
\end{enumerate}

С учетом задания на курсовой проект и выявленной специфики предметной области
основные функциональные требования к системе представлены в таблице~\ref{tab:functional_requirements}.
\clearpage

\begin{table}[ht!]
    \centering
    \small
    \caption{Функциональные требования к разрабатываемому веб-приложению}
    \label{tab:functional_requirements}
    \begin{tabular}{|p{0.11\textwidth}|p{0.54\textwidth}|p{0.21\textwidth}|}
        \hline
        \textbf{Код} & \textbf{Содержание требования} & \textbf{Назначение требования} \\ \hline
        FR-1 & Система должна поддерживать аутентификацию пользователей и последующую авторизацию при обращении к функциям веб-приложения. & Контролируемый вход в систему \\ \hline
        FR-2 & Система должна поддерживать ролевую модель доступа не менее чем для администратора, инженера и наблюдателя, а также ограничение доступа к данным по разрешенным объектам эксплуатации. & Безопасное разграничение прав \\ \hline
        FR-3 & Система должна обеспечивать CRUD-операции (Create, Read, Update, Delete --- создание, чтение, изменение и удаление) для сущностей \textit{Объект}, \textit{Контроллер}, \textit{Конфигурация}, \textit{Диагностическая запись}. & Выполнение базовых операций с данными \\ \hline
        FR-4 & Система должна поддерживать загрузку, скачивание и удаление файлов конфигураций, журналов, отчетов и других вложений, привязанных к сущностям системы. & Работа с инженерными файлами \\ \hline
        FR-5 & Система должна обеспечивать прием данных от внешнего локального клиента через серверный интерфейс прикладного программирования (API), сохраняя информацию об источнике и времени синхронизации. & Интеграция с локальным сбором данных \\ \hline
        FR-6 & Система должна предоставлять поиск, фильтрацию и агрегирование диагностических данных по объекту, контроллеру, типу события и периоду времени. & Поддержка анализа эксплуатации \\ \hline
        FR-7 & Система должна фиксировать автора изменения, время операции, комментарий и историю поступления конфигурационных данных. & Аудит и прослеживаемость \\ \hline
        FR-8 & Система не должна предоставлять в базовой версии функции удаленной записи конфигурации, перезапуска или иного управляющего воздействия на ПЛК через веб-интерфейс. & Снижение рисков для промышленной инфраструктуры \\ \hline
    \end{tabular}
    \normalsize
\end{table}

В качестве базовой предметной модели целесообразно использовать следующие
сущности: \textit{Пользователь}, \textit{Роль}, \textit{Объект эксплуатации},
\textit{Контроллер}, \textit{Конфигурация}, \textit{Диагностическая запись},
\textit{Файл-вложение}, \textit{Сессия синхронизации}. Такая модель сохраняет
необходимый акцент на веб-приложении, поскольку именно серверная система хранит
структурированные данные, проверяет права, предоставляет интерфейсы и формирует
аналитику. Локальный клиент в данном контуре является только поставщиком данных,
а не самостоятельным центром управления.
